\documentclass[11pt]{article}
\usepackage{fontspec}
\usepackage[french]{babel}
\usepackage{amsmath}
\usepackage{amssymb}
\usepackage{amsthm}
\usepackage{mathrsfs}
\usepackage{tikz-cd}
\usepackage{enumitem}

\newcommand{\printedpage}[1]{\marginpar{\footnotesize p.~#1}}
\newcommand{\msp}[1]{\marginpar{\footnotesize [#1]}}
\newcommand{\Ob}{\operatorname{Ob}}
\newcommand{\Ind}{\operatorname{Ind}}
\newcommand{\Filt}{\operatorname{Filt}}
\newcommand{\card}{\operatorname{card}}
\newcommand{\End}{\operatorname{End}}
\newcommand{\Fl}{\operatorname{Fl}}
\newcommand{\Kar}{\operatorname{Kar}}

\begin{document}

\printedpage{69}
% head: A. Grothendieck : Dérivateurs, Ch. XVIII, Édition des Universités de Montpellier II et Paris VII

\[
\begin{tikzcd}[column sep=small]
 & i(y) \arrow[dr, dotted, "p_\alpha q_y=p_x"] & \\
Z=i(y) \arrow[ur, "\mathrm{id}"] \arrow[dr, "q_y"'] & & X \ , \\
 & X_\alpha \arrow[ur, dotted, "p_\alpha"'] &
\end{tikzcd}
\]
on a donc, en posant $x=(y,p_\alpha q_y)\in\Ob C/X$,
\[
\underbrace{u_x\circ\mathrm{id}_{i(y)}}_{u_x} \ = \ u_\alpha q_y \ ,
\]
donc pour prouver $(**)$, il reste à vérifier
\[
u_x \ = \ \underbrace{(up_\alpha)q_y}_{u(p_\alpha q_y)\,=\,up_x} \ ,
\]
\msp{page 197}
qui est en effet vérifiée par définition de $u_x$. Cela achève de prouver $2^\circ$.

Prouvons $3^\circ$. Comme $i'$ est pleinement fidèle, en vertu de $1^\circ$, pour prouver que $i'$ est strictement génératrice, il suffit d'établir que pour tout objet $X$ de $\mathcal{M}$, la famille des flèches $i'(x')\longrightarrow X$, indexée par $\Ob C'/X$, est épimorphique stricte. Or cette famille contient la famille des $\underbrace{i'(\alpha(x))}_{=\,i(x)}\longrightarrow X$, indexée par $\Ob C/X$. Donc il résulte de l'hypothèse sur $i$, et de $2^\circ$, que cette famille est épimorphique stricte. On a gagné ! Ouf \dots

\textsc{Démonstration de 4.8.1.} On a
\[
\mathrm{c}' \ \Longrightarrow \ \mathrm{c} \ \Longrightarrow \ \mathrm{b} \ \Longrightarrow \ \mathrm{a} \, ,
\]
la première implication tautologique, la seconde résultant du fait que $\pi_0$ cardinal régulier pour la fcga envisagée, implique que
\[
\mathcal{M} \ \xleftarrow{\ \sim\ } \ \Ind_{\pi_0}(\mathcal{F}^{\pi_0}) \ ,
\]
$\mathcal{F}^{\pi_0}$ équivalente à une petite catégorie, d'où b).

\msp{page 198}
Enfin b) $\Longrightarrow$ a) résulte du fait que $\Ind_{\pi_0}(C)$ est stable par $\varinjlim$ grandes devant $\pi_0$ (4.4.15), et que $C$ est une sous-catégorie strictement génératrice de $\Ind_{\pi_0}(C)$, formée d'objets $\pi_0$-accessibles. Il nous faut expliciter le

\textbf{Lemme 4.8.4.} \textit{Soit $C$ une $\mathfrak{U}$-catégorie, $\pi_0$ un cardinal infini. Alors $C$ est strictement génératrice dans $\mathcal{M}=\Ind C$ et dans $\mathcal{M}=\Ind_{\pi_0}C$.}

\msp{page 199}
En effet, le foncteur canonique $\mathcal{M}\longrightarrow C^\wedge$ est pleinement fidèle, et on applique le critère c) de 4.8.2.

Pour achever de prouver 4.8.1, il reste donc à établir l'implication
\[
\mathrm{a} \ \Longrightarrow \ \mathrm{c}' \, .
\]

\printedpage{70}
% head: A. Grothendieck : Dérivateurs, Ch. XVIII, Édition des Universités de Montpellier II et Paris VII

Supposons donc que $\mathcal{M}$ satisfasse à $\mathrm{L}_{\pi_0}$, et admette une petite sous-catégorie strictement génératrice $S$, formée d'objets $\pi_0$-accessibles, donc un foncteur canonique
\[
\Ind_{\pi_0}S \ \longrightarrow \ \mathcal{M} \ ,
\]
prolongeant l'inclusion $S\hookrightarrow\mathcal{M}$, de façon à commuter aux $\varinjlim$ grandes devant $\pi_0$, et du fait que $S\subset\mathcal{M}_{\pi_0}$, ce foncteur est \textit{pleinement fidèle.} Mais comment prouver qu'il est essentiellement surjectif ? Ce ne sera d'ailleurs pas vrai sans autre hypothèse sur $S$. L'idée naturelle est de saturer $S$ pour les $\varinjlim$
\msp{page 200}
de cardinal $\leqslant\pi_0$ représentables dans $\mathcal{M}$, et notamment par facteurs directs, mais même alors (et même en admettant que $\mathcal{M}_{\pi_0}$ lui-même est équivalent à une petite catégorie, et prenant $C=\mathcal{M}_{\pi_0}$), il n'est pourtant pas clair que le foncteur pleinement fidèle $\Ind_{\pi_0}C\longrightarrow\mathcal{M}$ soit essentiellement surjectif, ce qui équivaudrait aussi à dire que pour tout $F$ dans $\mathcal{M}$, $C/F$ est grande devant $\pi_0$. Il n'est en effet pas même clair que $C/F$ soit filtrante.

Donc finalement, je doute qu'on ait a) $\Longrightarrow$ c$'$) ou seulement a) $\Longrightarrow$ b), la situation n'est pas bien comprise visiblement.

Mais prouvons quand même
\[
\mathrm{b} \ \Longrightarrow \ \mathrm{c}' \, .
\]
Pour ceci, on peut supposer $\mathcal{M}=\Ind_{\pi_0}C$, $C$ petite. On sait alors que la filtration de $\mathcal{M}$ par les $(\mathcal{M}_\pi)_{\pi\geqslant\pi_0}$ est une fcga (prop. 4.7.7 et définition 4.7.8), et même une fcg \textit{canonique.} De plus, je dis que le cardinal $\pi_0$ est \textit{utile} et même régulier pour cette filtration (ce qui achèvera de prouver c$'$). On a un résultat un peu plus précis :

\msp{page 201}
\textbf{Corollaire 4.8.5.} \textit{Soient $C$ une catégorie équivalente à une petite catégorie, $\pi_0$ un cardinal infini, $\mathcal{M}=\Ind_{\pi_0}(C)\subset C^\wedge$. Alors :}
\begin{enumerate}
\item[\textit{1$^\circ$)}] $\pi_0$ \textit{est un cardinal} utile \textit{et même régulier pour la fcg canonique} $(\mathcal{M}_\pi)_{\pi\geqslant\pi_0}$ \textit{de} $\mathcal{M}$.
\item[\textit{2$^\circ$)}] $\mathcal{M}_{\pi_0}=C^\natural$ \textit{(sous-catégorie pleine de $\mathcal{M}=\Ind_{\pi_0}C$ formée des facteurs directs dans $\mathcal{M}$ d'objets de $C$), et $\mathcal{M}_{\pi_0}$ est une enveloppe de Karoubi de $C$.}
\end{enumerate}
(\footnote{\textbf{Corollaire 4.8.5.1.} \textit{Si $C$ est karoubienne, alors $\mathcal{M}_{\pi_0}\subset C^\wedge$ est la catégorie formée des foncteurs} représentables \textit{dans} $C\overset{\sim}{\hookrightarrow}\mathcal{M}_{\pi_0}$.}, \footnote{\textbf{4.8.5.2.} Expliciter aussi ceci : Soit $u : C\longrightarrow C'$ un foncteur, avec $C$, $C'$ équivalentes à des petites catégories. Conditions équivalentes :
\begin{enumerate}
\item[a)] $u$ induit $\Kar C\xrightarrow{\ \sim\ }\Kar C'$.
\item[a$'$)] $u$ pleinement fidèle, et tout objet de $C'$ est isomorphe à un facteur direct d'un objet de l'image essentielle.
\item[b)] $\Ind(u) : \Ind(C)\longrightarrow\Ind(C')$ est une équivalence.
\item[c)] $\Ind_{\pi_0}(u) : \Ind_{\pi_0}(C)\longrightarrow\Ind_{\pi_0}(C')$ est une équivalence ($\pi_0$ un cardinal infini donné).
\item[d)] $u_! : C^\wedge\longrightarrow C'^\wedge$ est une équivalence.
\item[d$'$)] $u_* : C'^\wedge\longrightarrow C^\wedge$ est une équivalence \texttt{[plutôt} $u_* : C^\wedge\longrightarrow C'^\wedge$\texttt{]}.
\item[d$''$)] $u^* : C'^\wedge\longrightarrow C^\wedge$ est une équivalence.
\end{enumerate}}).

\printedpage{71}
% head: A. Grothendieck : Dérivateurs, Ch. XVIII, Édition des Universités de Montpellier II et Paris VII

\textsc{Démonstration.}

$1^\circ$) On sait déjà que
\[
\Ind_{\pi_0}\mathcal{M}_{\pi_0} \ \longrightarrow \ \mathcal{M} \tag*{$(*)$}
\]
est pleinement fidèle, et pour voir que c'est essentiellement surjectif, on note que l'on a
\[
\mathcal{M} \ \simeq \ \Ind_{\pi_0}C \ \hookrightarrow \ \Ind_{\pi_0}\mathcal{M}_{\pi_0} \ \longrightarrow \ \mathcal{M} \ ,
\]
le composé étant isomorphe au foncteur identique, est essentiellement surjectif, donc aussi $(*)$. Cela prouve donc que $\pi_0$ est un cardinal utile pour la filtration envisagée, prouvons que ce cardinal est même régulier, donc que pour
\[
F \ \in \ \Ob\mathcal{M}_\pi \ , \qquad \text{avec } \pi>\pi_0,\ \pi^{\pi_0}\geqslant\pi_1
\]
on a
\[
F \ = \ \varinjlim_I x_i \ , \qquad x_i\in\Ob\mathcal{M}_{\pi_0},\ I \text{ grande devant } \pi_0,\ \card I\leqslant\pi^{\pi_0}.
\]
\msp{page 202}
Pour ceci, on utilise (4.7.6.2), page 171,
\[
\mathcal{M}_\pi \ \subset \ \Filt^{\pi^{\pi_0}} \qquad (\text{si } \pi\geqslant\pi_0,\ \pi^{\pi_0}\geqslant\pi_1),
\]
et la définition de $\Filt^{\pi'}$ (pour $\pi'\geqslant\pi_0$) nous montre que c'est là la conclusion cherchée.

Prouvons donc $2^\circ$). Comme $\mathcal{M}_{\pi_0}$ est stable par les $\varinjlim$ inductives finies qui sont représentables dans $\mathcal{M}$ (SGA 4 I 9.9), il est stable par facteurs directs, et comme il contient $C$, on a donc $C^\natural\subset\mathcal{M}_{\pi_0}$. L'inclusion inverse $\mathcal{M}_{\pi_0}\subset C^\natural$ se voit de la façon standard, en écrivant un objet $F$ de $\mathcal{M}_{\pi_0}$ comme $\varinjlim\limits_I x_i$, les $x_i$ dans $C$, $I$ grande devant $\pi_0$, alors $F$ est facteur direct d'un des $x_i$.

Le fait que $C^\natural$ est une \textit{enveloppe de Karoubi} de $C$ dans $\mathcal{M}\subset C^\wedge$ provient du fait que $\mathcal{M}$ est stable par facteurs directs, et que $C$ est une sous-catégorie pleine de $\mathcal{M}$, cf. SGA 4 I (\,?\,\dots).

\textbf{4.8.6. Questions ouvertes.} On a ainsi prouvé 4.8.1 revu et corrigé. Il reste la question :
\begin{enumerate}
\item[1)] Donner un exemple où 4.8.1 a) est vérifié, mais non b) etc. (\footnote{Il suffit de trouver un cardinal non régulier, pour la fcg canonique d'une $\mathcal{M}=\Ind_{\pi_0}C$.}).
\item[2)] Dégager des conditions supplémentaires serviables qui, jointes à a), impliquent b), donc c, c$'$.
\end{enumerate}

\msp{page 203}
Signalons aussi pour mémoire les questions suivantes.
\begin{enumerate}
\item[3)] Soit $\mathcal{M}=\Ind_{\pi_0}(C)$, et $\pi>\pi_0$. Alors $\Filt^\pi$ est-il stable par facteurs directs ? Cela est-il vrai du moins si $\pi=\pi^{\pi_0}$, ou ce qui revient au même (cf. 4.7.3), a-t-on dans ce cas $\Filt^\pi=\Filt'^{\pi}$ ? Si oui, il serait inutile dans la description 4.7.7 des filtrations cardinales, dans la majoration de cardinaux (4.7.7.1), $\card I\leqslant\pi'^{\pi}$, de mettre l'hypothèse restrictive $\pi'^{\pi}\geqslant\pi_1$.
\end{enumerate}

Pour avoir une réponse affirmative à cette question, pour $\pi$ donné, il suffirait (et probablement faudrait) d'avoir une réponse affirmative à celle-ci : Soient $\mathcal{D}$, $\mathcal{D}_0$ deux catégories grandes devant $\pi_0$, $\mathcal{D}_0$ facteur direct de $\mathcal{D}$, $\mathcal{D}$ admettant une catégorie cofinale $I$ grande devant $\pi_0$, de cardinal $\leqslant\pi$. Alors en est-il de même de $\mathcal{D}_0$ ? Si $p$ est une rétraction de $\mathcal{D}$ sur $\mathcal{D}_0$, la difficulté est que je ne sais pas si $p$ est foncteur cofinal. En fait, je vois que
\printedpage{72}
% head: A. Grothendieck : Dérivateurs, Ch. XVIII, Édition des Universités de Montpellier II et Paris VII
\msp{page 204}
ce n'est pas toujours le cas, même si $\mathcal{D}$ est finie. Prenons p.ex. pour $\mathcal{D}_0$ une catégorie ayant objet final $e$, et une flèche $x\longrightarrow e$ qui n'est pas un isomorphisme, et admettant une section (p.ex. $\mathcal{D}_0$ la sous-catégorie pleine de $\boldsymbol{\Delta}$ formée par $\boldsymbol{\Delta}_0$, $\boldsymbol{\Delta}_1$ : \begin{tikzcd}[cramped] \boldsymbol{\Delta}_0 \arrow[r, shift left, "\partial_0"] \arrow[r, shift right, "\partial_1"'] & \boldsymbol{\Delta}_1 \arrow[l, bend left=60, "\kappa"] \end{tikzcd} ). Soit $\mathcal{D}$ déduit de $\mathcal{D}_0$ en ajoutant un objet final $\tilde{e}$. Alors $\mathcal{D}_0$, $\mathcal{D}$ sont filtrantes, et on obtient une rétraction de $\mathcal{D}$ sur $\mathcal{D}_0$, envoyant $\tilde{e}$ sur $x$. Ce foncteur n'est pas cofinal, car un foncteur cofinal doit envoyer $\tilde{e}$ en $e$.

Plus généralement, si $\mathcal{D}_0$ est une catégorie telle qu'il existe une flèche
\[
\alpha : e_{\mathcal{D}_0^\wedge} \ \longrightarrow \ x_0 \tag*{$(*)$}
\]
dans $\mathcal{D}_0^\wedge$, avec $x_0\in\Ob\mathcal{D}_0$, alors on voit que $\mathcal{D}_0$ est grande devant $\pi_0$ quel que soit le cardinal infini $\pi_0$, et que d'autre part, désignant encore par $\mathcal{D}$ la catégorie déduite de $\mathcal{D}_0$ en lui ajoutant un objet final, il y a une rétraction de $\mathcal{D}$ sur $\mathcal{D}_0$ qui envoie $e_{\mathcal{D}}$ dans $x_0$. Cette rétraction n'est cofinale que si $x_0$ est un objet final de $\mathcal{D}_0$, i.e. si et seulement si $(*)$ est un isomorphisme. Supposons qu'il n'en soit pas ainsi. Mais il y a un foncteur cofinal de la sous-catégorie $I_0$ de $\mathcal{D}_0$, telle que $\Ob I_0=\{x_0\}$, $\Fl I_0=\End_{I_0}(x_0)=\{p,\mathrm{id}_{x_0}\}$, où $p$ est le composé $x_0\xrightarrow{\text{can.}}e_{\mathcal{D}_0^\wedge}\xrightarrow{\alpha}x_0$. Donc pas de contre-exemple de cette façon !

\msp{page 205}
\begin{enumerate}
\item[4)] Considérons $\mathcal{M}=\Ind_{\pi_0}(C)$, $C$ petite, et considérons la filtration canonique $(\mathcal{M}_\pi)_{\pi\geqslant\pi_0}$. Donner des exemples où il y a des cardinaux non utiles, ou mieux, ou pour tout cardinal $\pi$, existe un cardinal $\pi'>\pi$ qui soit non utile. (Cela donnera l'exemple demandé dans $1^\circ$.) D'autre part, donner des exemples de cardinaux utiles qui ne soient réguliers -- si possible, comme tantôt, de tels cardinaux aussi grands qu'on veut.
\end{enumerate}

\textbf{Définition 4.8.7.} Soit $\mathcal{M}$ une $\mathfrak{U}$-catégorie, $\pi_0$ un cardinal infini. On dit que $\mathcal{M}$ est $\pi_0$-\textit{accessible} si $\mathcal{M}$ est équivalente à une catégorie $\Ind_{\pi_0}(C)$, $C$ une petite catégorie, i.e. si elle satisfait aux conditions équivalentes b, c, c$'$ de 4.8.1 (impliquant la condition a), sans doute strictement plus faible). On dit que $\mathcal{M}$ est accessible s'il existe un cardinal tel que $\mathcal{M}$ soit $\pi_0$-accessible (\textbf{N.B.} si $\pi_0$ est un tel cardinal, alors tout cardinal $\pi>\pi_0$ tel que $\pi^{\pi_0}=\pi$, notamment les cardinaux $\pi'^{\pi_0}$, avec $\pi'\geqslant 2$, ont la même propriété), ou ce
\msp{page 206}
qui revient au même d'après 4.8.1, si $\mathcal{M}$ admet une fcga, ou encore une fcg canonique.

\textbf{Commentaire 4.8.8.} Cette définition est en contradiction avec la terminologie introduite dans SGA 4 I 9.4 (p. 114) ; on désignera dorénavant sous les noms de ``catégorie faiblement $\pi_0$-accessible'', ``catégorie faiblement accessible'' les notions introduites dans 9.4 -- i.e. les catégories admettant une sous-catégorie pleine génératrice (pas même strictement génératrice) dont les objets sont $\pi_0$-accessibles (resp. accessibles).

Notons que si $\mathcal{M}$ est $\pi_0$-accessible (resp. accessible), la condition de 4.8.1 nous dit que $\mathcal{M}$ est faiblement $\pi_0$-accessible (resp. faiblement accessible), et même qu'il existe une petite sous-catégorie pleine \textit{strictement} génératrice, qui soit formée d'objets $\pi_0$-accessibles (resp. accessibles). Je doute fort que cette dernière condition nécessaire soit aussi suffisante pour
\printedpage{73}
% head: A. Grothendieck : Dérivateurs, Ch. XVIII, Édition des Universités de Montpellier II et Paris VII
que $\mathcal{M}$ soit $\pi_0$-accessible (resp. accessible). Mais en vertu de 4.6.3.3 (p. 150), il en est ainsi si on suppose $C$ stable par petites $\varinjlim$ , ou ce qui revient au même, moyennant $\mathrm{L}_{\pi_0}$, stable par limites inductives de cardinal $\leqslant\pi_0$.

\msp{page 207}
\textbf{Proposition 4.8.9.} \textit{Soit $\mathcal{M}$ une $\mathfrak{U}$-catégorie accessible (4.8.7), $\pi_0$ un cardinal tel que $\mathcal{M}$ satisfait à $L_{\pi_0}$. Conditions équivalentes :}
\begin{enumerate}
\item[\textit{a)}] \textit{$\mathcal{M}$ est une catégorie paratopos (cf. 4.6.3, $3^\circ$, p. 150).}
\item[\textit{b)}] \textit{$\mathcal{M}$ est stable par petites limites inductives.}
\item[\textit{b$'$)}] \textit{$\mathcal{M}$ est stable par limites inductives de cardinal $\leqslant\pi_0$.}
\item[\textit{b$''$)}] \textit{(Si $\mathcal{M}$ est équivalente à une catégories} $\Ind_{\pi_0}C$\textit{, avec $C$ une petite catégorie} karoubienne.\textit{) La catégorie $C$ est stable (en elle-même) par limites inductives de cardinal $\leqslant\pi_0$.} % [sic] une catégories -> une catégorie, accord du nombre
\item[\textit{c)}] \textit{$\mathcal{M}$ est stable par petites limites projectives.}
\item[\textit{c$'$)}] \textit{(Sous l'hypothèse $\mathcal{M}\simeq\Ind_{\pi_0}(C)$, et quand on se donne un cardinal $\pi$ tel que $\pi=\pi^{\pi_0}\geqslant\pi_1=\sup(\pi_0,\card\textup{réd}\,C)$.) $\mathcal{M}_\pi$ est stable par $\varprojlim$ finies, et $\mathcal{M}$ est stable par petits produits infinis (ou encore, par $\varprojlim$ filtrantes).}
\item[\textit{c$''$)}] \textit{Il existe un cardinal infini $\pi$, et une petite catégorie $\mathcal{D}$ stable par $\varprojlim$ finies, telle que $\mathcal{M}\simeq\Ind_\pi(\mathcal{D})$, de plus $\mathcal{M}$ stable par produits infinis (ou encore par $\varprojlim$ filtrantes).}
\end{enumerate}
(\footnote{\textbf{N.B.} L'équivalence de c$'$, c$''$ avec les autres conditions n'est pas établie (ici, mais seulement dans 4.11 plus bas.)}, \footnote{\textbf{N.B.} Les critères c$'$, c$''$ me paraissent plus techniques et moins utiles \texttt{[c$'$, c$''$ ci-dessus sont mis en parenthèses]}.}).

\textsc{Démonstration.} Les implications suivantes sont tautologiques, ou déjà $\pm$ connues :
\[
\begin{tikzcd}[column sep=large, row sep=large]
\mathrm{a} \arrow[r, Rightarrow] & \mathrm{b} \arrow[r, Rightarrow, "\text{cf. 4.4.14}"] \arrow[dl, Rightarrow] & \mathrm{b}' \arrow[r, Rightarrow, "\bullet"] & \mathrm{b}'' \\
\mathrm{c}\ ( \arrow[r, Rightarrow, "\bullet"] & \mathrm{c}' \arrow[r, Rightarrow] & \mathrm{c}'')\ , &
\end{tikzcd}
\]
\msp{page 208}
où seules les implications pointées b$'\Longrightarrow$ b$''$ et c $\Longrightarrow$ c$'$ demandent une démonstration. Pour b$'\Longrightarrow$ b$''$, on note que l'hypothèse $C$ karoubienne implique que $C\simeq\mathcal{M}_{\pi_0}$ (4.8.5, $2^\circ$), or on sait que $\mathcal{M}_{\pi_0}$ est stable dans $\mathcal{M}$ par limites inductives de cardinal $\leqslant\pi_0$ (et celles-ci sont représentables dans $\mathcal{M}$ en vertu de b$'$), donc $\mathcal{M}_{\pi_0}$ admet des $\varinjlim$ du type envisagé, donc $C$ aussi.

c $\Longrightarrow$ c$'$. Notons que si $\mathcal{M}=\Ind_{\pi_0}(C)$, alors les $\varprojlim$ représentables dans $\mathcal{M}$ se calculent dans $C^\wedge$. Il s'ensuit que si $\pi\geqslant\pi_1$, alors $\Filt^\pi$ est stable par $\varprojlim$ finies ($\Filt^\pi$ étant donné alors par la condition a$''$ de 4.7.1), d'autre part, si de plus $\pi^{\pi_0}=\pi$, $\pi$ est un cardinal utile, donc $\mathcal{M}\simeq\Ind_\pi(\mathcal{D})$, où $\mathcal{D}=\Filt^\pi=\mathcal{M}_\pi$. Il suffit maintenant de voir que $\Ind_\pi(\mathcal{D})$ est alors stable par $\varprojlim$ finies, pour conclure qu'il en est de même de $\mathcal{M}$, donc que $\mathcal{M}$ est stable par limites projectives quelconques, s'il l'est par petites produites ou par petites $\varprojlim$ filtrantes, d'où c$'$. Il reste donc à prouver le lemme suivant. % [sic] petites produites -> petits produits, accord

\end{document}
